\documentclass[10pt,a4paper]{article}

%double si on veut une version prof et une version élève. simple si on veut une seule version.
\newcommand{\buildmode}{simple}

\InputIfFileExists{version.tex}{}{}
% Version (définie par le script)
\providecommand{\version}{eleve}

\usepackage{feuilleinterro}

%%%à modifier à chaque fois

\renewcommand{\niveau}{SBM 2}
\renewcommand{\chapitre}{Fonctions, généralités}
\renewcommand{\numchap}{0.2}
\renewcommand{\theme}{Fonctions}

\begin{document}

% =========================================================
% PAGE 1 : VERSIONS A ET B
% =========================================================


\interroversion{A}

\textbf{Exercice 1 -- Tableau de valeurs \hfill /2}

On considère la fonction \(f\) définie  par \(f(x)=x^2-2x-3\).

Complétez le tableau de valeurs suivant (une valeur est déjà calculée à titre d'exemple : \(f(1)=1^2-2\times 1-3=-4\)).

\begin{center}
    \renewcommand{\arraystretch}{1.5}
    \setlength{\tabcolsep}{0.25cm}
    \begin{tabular}{|*{7}{c|}}
    \hline
    \(x\) & \(-2\) & \(-1\) & \(0\) & \(1\) & \(2\) & \(3\) \\
    \hline
    \(f(x)\) & & \(0\) & & \(-4\) & \(-3\) & \(0\) \\
    \hline
    \end{tabular}
\end{center}


\textbf{Exercice 2 -- Signe d'un produit de fonctions affines \hfill /4}

On considère les fonctions affines \(f\) et \(g\) définies par
\[
f(x)=2x-3 \qquad\text{et}\qquad g(x)=-x+4,
\]
ainsi que la fonction \(h\) définie  par \(h(x)=(2x-3)(-x+4)\).

\begin{enumerate}
\item Faire les tableaux de signes des fonctions \(f\) et \(g\)

\vspace{6cm}

\item En déduire le tableau de signes de \(h\) .

\vspace{4cm}

\end{enumerate}

\text{ }
\newpage

\interroversion{B}

\textbf{Exercice 1 -- Tableau de valeurs \hfill /2}

On considère la fonction \(f\) définie sur \(\mathbb{R}\) par \(f(x)=x^2+x-2\).

Complétez le tableau de valeurs suivant (une valeur est déjà calculée à titre d'exemple : \(f(2)=2^2+2-2=4\)).

\begin{center}
    \renewcommand{\arraystretch}{1.5}
    \setlength{\tabcolsep}{0.25cm}
    \begin{tabular}{|*{7}{c|}}
    \hline
    \(x\) & \(-2\) & \(-1\) & \(0\) & \(1\) & \(2\) & \(3\) \\
    \hline
    \(f(x)\) & \(0\) & & \(-2\) & \(0\) & \(4\) & \\
    \hline
    \end{tabular}
\end{center}

\textbf{Exercice 2 -- Signe d'un produit de fonctions affines \hfill /4}

On considère les fonctions affines \(f\) et \(g\) définies sur par
\[
f(x)=-3x+6 \qquad\text{et}\qquad g(x)=x+1,
\]
ainsi que la fonction \(h\) définie sur par \(h(x)=(-3x+6)(x+1)\).


\begin{enumerate}
\item Faire les tableaux de signes des fonctions \(f\) et \(g\)

\vspace{6cm}

\item En déduire le tableau de signes de \(h\) .

\vspace{4cm}

\end{enumerate}

\text{ }

\newpage

% =========================================================
% PAGE 2 : VERSIONS C ET D
% =========================================================

\interroversion{C}

\textbf{Exercice 1 -- Tableau de valeurs \hfill /2}

On considère la fonction \(f\) définie sur \(\mathbb{R}\) par \(f(x)=x^2-3x\).

Complétez le tableau de valeurs suivant (une valeur est déjà calculée à titre d'exemple : \(f(-1)=(-1)^2-3\times(-1)=4\)).

\begin{center}
    \renewcommand{\arraystretch}{1.5}
    \setlength{\tabcolsep}{0.25cm}
    \begin{tabular}{|*{7}{c|}}
    \hline
    \(x\) & \(-2\) & \(-1\) & \(0\) & \(1\) & \(2\) & \(3\) \\
    \hline
    \(f(x)\) & & \(4\) & \(0\) & & \(-2\) & \(0\) \\
    \hline
    \end{tabular}
\end{center}

\textbf{Exercice 2 -- Signe d'un produit de fonctions affines \hfill /4}

On considère les fonctions affines \(f\) et \(g\) définies par
\[
f(x)=4x-2 \qquad\text{et}\qquad g(x)=-2x-4,
\]
ainsi que la fonction \(h\) définie  par \(h(x)=(4x-2)(-2x-4)\).

\begin{enumerate}
\item Faire les tableaux de signes des fonctions \(f\) et \(g\)

\vspace{6cm}

\item En déduire le tableau de signes de \(h\) .

\vspace{4cm}

\end{enumerate}

\text{ }
\newpage

\interroversion{D}

\textbf{Exercice 1 -- Tableau de valeurs \hfill /2}

On considère la fonction \(f\) définie sur \(\mathbb{R}\) par \(f(x)=x^2+2x-1\).

Complétez le tableau de valeurs suivant (une valeur est déjà calculée à titre d'exemple : \(f(0)=0^2+2\times 0-1=-1\)).

\begin{center}
    \renewcommand{\arraystretch}{1.5}
    \setlength{\tabcolsep}{0.25cm}
    \begin{tabular}{|*{7}{c|}}
    \hline
    \(x\) & \(-2\) & \(-1\) & \(0\) & \(1\) & \(2\) & \(3\) \\
    \hline
    \(f(x)\) & \(-1\) & & \(-1\) & \(2\) & & \(14\) \\
    \hline
    \end{tabular}
\end{center}

\textbf{Exercice 2 -- Signe d'un produit de fonctions affines \hfill /4}

On considère les fonctions affines \(f\) et \(g\) définies par
\[
f(x)=-x+3 \qquad\text{et}\qquad g(x)=3x+6,
\]
ainsi que la fonction \(h\) définie par \(h(x)=(-x+3)(3x+6)\).


\begin{enumerate}
\item Faire les tableaux de signes des fonctions \(f\) et \(g\)

\vspace{6cm}

\item En déduire le tableau de signes de \(h\) .

\vspace{4cm}

\end{enumerate}

\text{ }

\end{document}
